% Options for packages loaded elsewhere
\PassOptionsToPackage{unicode}{hyperref}
\PassOptionsToPackage{hyphens}{url}
\documentclass[
]{article}
\usepackage{xcolor}
\usepackage[margin=1in]{geometry}
\usepackage{amsmath,amssymb}
\setcounter{secnumdepth}{-\maxdimen} % remove section numbering
\usepackage{iftex}
\ifPDFTeX
  \usepackage[T1]{fontenc}
  \usepackage[utf8]{inputenc}
  \usepackage{textcomp} % provide euro and other symbols
\else % if luatex or xetex
  \usepackage{unicode-math} % this also loads fontspec
  \defaultfontfeatures{Scale=MatchLowercase}
  \defaultfontfeatures[\rmfamily]{Ligatures=TeX,Scale=1}
\fi
\usepackage{lmodern}
\ifPDFTeX\else
  % xetex/luatex font selection
\fi
% Use upquote if available, for straight quotes in verbatim environments
\IfFileExists{upquote.sty}{\usepackage{upquote}}{}
\IfFileExists{microtype.sty}{% use microtype if available
  \usepackage[]{microtype}
  \UseMicrotypeSet[protrusion]{basicmath} % disable protrusion for tt fonts
}{}
\makeatletter
\@ifundefined{KOMAClassName}{% if non-KOMA class
  \IfFileExists{parskip.sty}{%
    \usepackage{parskip}
  }{% else
    \setlength{\parindent}{0pt}
    \setlength{\parskip}{6pt plus 2pt minus 1pt}}
}{% if KOMA class
  \KOMAoptions{parskip=half}}
\makeatother
\usepackage{color}
\usepackage{fancyvrb}
\newcommand{\VerbBar}{|}
\newcommand{\VERB}{\Verb[commandchars=\\\{\}]}
\DefineVerbatimEnvironment{Highlighting}{Verbatim}{commandchars=\\\{\}}
% Add ',fontsize=\small' for more characters per line
\usepackage{framed}
\definecolor{shadecolor}{RGB}{248,248,248}
\newenvironment{Shaded}{\begin{snugshade}}{\end{snugshade}}
\newcommand{\AlertTok}[1]{\textcolor[rgb]{0.94,0.16,0.16}{#1}}
\newcommand{\AnnotationTok}[1]{\textcolor[rgb]{0.56,0.35,0.01}{\textbf{\textit{#1}}}}
\newcommand{\AttributeTok}[1]{\textcolor[rgb]{0.13,0.29,0.53}{#1}}
\newcommand{\BaseNTok}[1]{\textcolor[rgb]{0.00,0.00,0.81}{#1}}
\newcommand{\BuiltInTok}[1]{#1}
\newcommand{\CharTok}[1]{\textcolor[rgb]{0.31,0.60,0.02}{#1}}
\newcommand{\CommentTok}[1]{\textcolor[rgb]{0.56,0.35,0.01}{\textit{#1}}}
\newcommand{\CommentVarTok}[1]{\textcolor[rgb]{0.56,0.35,0.01}{\textbf{\textit{#1}}}}
\newcommand{\ConstantTok}[1]{\textcolor[rgb]{0.56,0.35,0.01}{#1}}
\newcommand{\ControlFlowTok}[1]{\textcolor[rgb]{0.13,0.29,0.53}{\textbf{#1}}}
\newcommand{\DataTypeTok}[1]{\textcolor[rgb]{0.13,0.29,0.53}{#1}}
\newcommand{\DecValTok}[1]{\textcolor[rgb]{0.00,0.00,0.81}{#1}}
\newcommand{\DocumentationTok}[1]{\textcolor[rgb]{0.56,0.35,0.01}{\textbf{\textit{#1}}}}
\newcommand{\ErrorTok}[1]{\textcolor[rgb]{0.64,0.00,0.00}{\textbf{#1}}}
\newcommand{\ExtensionTok}[1]{#1}
\newcommand{\FloatTok}[1]{\textcolor[rgb]{0.00,0.00,0.81}{#1}}
\newcommand{\FunctionTok}[1]{\textcolor[rgb]{0.13,0.29,0.53}{\textbf{#1}}}
\newcommand{\ImportTok}[1]{#1}
\newcommand{\InformationTok}[1]{\textcolor[rgb]{0.56,0.35,0.01}{\textbf{\textit{#1}}}}
\newcommand{\KeywordTok}[1]{\textcolor[rgb]{0.13,0.29,0.53}{\textbf{#1}}}
\newcommand{\NormalTok}[1]{#1}
\newcommand{\OperatorTok}[1]{\textcolor[rgb]{0.81,0.36,0.00}{\textbf{#1}}}
\newcommand{\OtherTok}[1]{\textcolor[rgb]{0.56,0.35,0.01}{#1}}
\newcommand{\PreprocessorTok}[1]{\textcolor[rgb]{0.56,0.35,0.01}{\textit{#1}}}
\newcommand{\RegionMarkerTok}[1]{#1}
\newcommand{\SpecialCharTok}[1]{\textcolor[rgb]{0.81,0.36,0.00}{\textbf{#1}}}
\newcommand{\SpecialStringTok}[1]{\textcolor[rgb]{0.31,0.60,0.02}{#1}}
\newcommand{\StringTok}[1]{\textcolor[rgb]{0.31,0.60,0.02}{#1}}
\newcommand{\VariableTok}[1]{\textcolor[rgb]{0.00,0.00,0.00}{#1}}
\newcommand{\VerbatimStringTok}[1]{\textcolor[rgb]{0.31,0.60,0.02}{#1}}
\newcommand{\WarningTok}[1]{\textcolor[rgb]{0.56,0.35,0.01}{\textbf{\textit{#1}}}}
\usepackage{graphicx}
\makeatletter
\newsavebox\pandoc@box
\newcommand*\pandocbounded[1]{% scales image to fit in text height/width
  \sbox\pandoc@box{#1}%
  \Gscale@div\@tempa{\textheight}{\dimexpr\ht\pandoc@box+\dp\pandoc@box\relax}%
  \Gscale@div\@tempb{\linewidth}{\wd\pandoc@box}%
  \ifdim\@tempb\p@<\@tempa\p@\let\@tempa\@tempb\fi% select the smaller of both
  \ifdim\@tempa\p@<\p@\scalebox{\@tempa}{\usebox\pandoc@box}%
  \else\usebox{\pandoc@box}%
  \fi%
}
% Set default figure placement to htbp
\def\fps@figure{htbp}
\makeatother
\setlength{\emergencystretch}{3em} % prevent overfull lines
\providecommand{\tightlist}{%
  \setlength{\itemsep}{0pt}\setlength{\parskip}{0pt}}
\usepackage{bookmark}
\IfFileExists{xurl.sty}{\usepackage{xurl}}{} % add URL line breaks if available
\urlstyle{same}
\hypersetup{
  pdftitle={Jeffreys' Prior},
  hidelinks,
  pdfcreator={LaTeX via pandoc}}

\title{Jeffreys' Prior}
\author{}
\date{\vspace{-2.5em}2026-04-17}

\begin{document}
\maketitle

\subsection{Introduction}\label{introduction}

Jeffreys' prior is the classical ``objective'' or ``reference'' prior in
Bayesian inference. Defined as proportional to the square root of the
Fisher information, it has the elegant property of being invariant under
monotone reparameterisation: the prior on a transformation
\(\eta = g(\theta)\) derived from the Jeffreys prior on \(\theta\)
equals the Jeffreys prior derived directly for \(\eta\). This invariance
--- absent from a flat prior, which is parameterisation-dependent --- is
what historically made Jeffreys' rule the default attempt at a
non-informative prior. In modern Bayesian practice, however, weakly
informative priors have largely displaced it because they regularise
small-sample inference more reliably and behave better in hierarchical
models.

\subsection{Prerequisites}\label{prerequisites}

A working understanding of Fisher information, the role of
reparameterisation in inference, and the difference between proper and
improper priors.

\subsection{Theory}\label{theory}

For a scalar parameter \(\theta\) with likelihood \(p(y \mid \theta)\),
the Jeffreys prior is

\[p(\theta) \propto \sqrt{I(\theta)}, \qquad I(\theta) = -\mathbb E\!\left[\frac{\partial^2 \log p(y \mid \theta)}{\partial \theta^2}\right].\]

Standard examples:

\begin{itemize}
\tightlist
\item
  Binomial\((n, \theta)\): \(I(\theta) \propto 1 / [\theta(1-\theta)]\),
  so the Jeffreys prior is
  \(\mathrm{Beta}(\tfrac{1}{2}, \tfrac{1}{2})\).
\item
  Poisson\((\lambda)\): \(I(\lambda) \propto 1 / \lambda\), so
  \(p(\lambda) \propto \lambda^{-1/2}\).
\item
  Normal mean (known variance): flat prior \(p(\mu) \propto 1\).
\item
  Normal scale (known mean): \(p(\sigma) \propto 1 / \sigma\).
\end{itemize}

Jeffreys priors are often improper, but the posterior is usually proper
as long as the sample size is sufficient. For multivariate parameters,
\(I(\theta)\) is a matrix and the prior is \(\sqrt{|I(\theta)|}\).

\subsection{Assumptions}\label{assumptions}

A regular likelihood (twice differentiable in \(\theta\), with finite
Fisher information across the parameter space). Boundary problems and
parameter-dimension growth complicate the construction; reference priors
(Bernardo) generalise Jeffreys' rule for hierarchical and
high-dimensional cases.

\subsection{R Implementation}\label{r-implementation}

\begin{Shaded}
\begin{Highlighting}[]
\CommentTok{\# Binomial: posterior under Jeffreys is Beta(1/2 + x, 1/2 + n {-} x)}
\NormalTok{n }\OtherTok{\textless{}{-}} \DecValTok{10}\NormalTok{; x }\OtherTok{\textless{}{-}} \DecValTok{8}
\NormalTok{post\_a }\OtherTok{\textless{}{-}} \FloatTok{0.5} \SpecialCharTok{+}\NormalTok{ x}
\NormalTok{post\_b }\OtherTok{\textless{}{-}} \FloatTok{0.5} \SpecialCharTok{+}\NormalTok{ n }\SpecialCharTok{{-}}\NormalTok{ x}
\FunctionTok{c}\NormalTok{(}\AttributeTok{mean =}\NormalTok{ post\_a }\SpecialCharTok{/}\NormalTok{ (post\_a }\SpecialCharTok{+}\NormalTok{ post\_b),}
  \AttributeTok{ci95 =} \FunctionTok{qbeta}\NormalTok{(}\FunctionTok{c}\NormalTok{(}\FloatTok{0.025}\NormalTok{, }\FloatTok{0.975}\NormalTok{), post\_a, post\_b))}

\CommentTok{\# Compare with uniform Beta(1, 1) and Beta(2, 2)}
\FunctionTok{data.frame}\NormalTok{(}
  \AttributeTok{prior =} \FunctionTok{c}\NormalTok{(}\StringTok{"Uniform(0,1)"}\NormalTok{, }\StringTok{"Jeffreys"}\NormalTok{, }\StringTok{"Beta(2,2)"}\NormalTok{),}
  \AttributeTok{post\_mean =} \FunctionTok{c}\NormalTok{((}\DecValTok{1} \SpecialCharTok{+}\NormalTok{ x)}\SpecialCharTok{/}\NormalTok{(}\DecValTok{2} \SpecialCharTok{+}\NormalTok{ n), (}\FloatTok{0.5} \SpecialCharTok{+}\NormalTok{ x)}\SpecialCharTok{/}\NormalTok{(}\DecValTok{1} \SpecialCharTok{+}\NormalTok{ n), (}\DecValTok{2} \SpecialCharTok{+}\NormalTok{ x)}\SpecialCharTok{/}\NormalTok{(}\DecValTok{4} \SpecialCharTok{+}\NormalTok{ n))}
\NormalTok{)}
\end{Highlighting}
\end{Shaded}

\subsection{Output \& Results}\label{output-results}

Closed-form summaries under three priors. With 8 successes in 10 trials
the Jeffreys posterior mean is approximately 0.77 (vs.~0.75 under the
uniform \(\mathrm{Beta}(1, 1)\) and 0.71 under the more informative
\(\mathrm{Beta}(2, 2)\)), and the 95 \% credible interval is roughly
\((0.47, 0.94)\). The comparison highlights how mildly Jeffreys differs
from a flat prior in this setting.

\subsection{Interpretation}\label{interpretation}

A reporting sentence: ``Jeffreys' prior \(\mathrm{Beta}(0.5, 0.5)\) on
the success probability combined with 8 successes in 10 trials yielded a
posterior mean of 0.77 (95 \% credible interval 0.47 to 0.94); the
result is essentially equivalent to a uniform prior at this sample size,
justifying the reference-prior framing.'' When Jeffreys differs
materially from a weakly informative prior, the prior is doing real work
and you should be explicit about why you chose it.

\subsection{Practical Tips}\label{practical-tips}

\begin{itemize}
\tightlist
\item
  The invariance property is mathematically attractive but in practice
  rarely affects applied conclusions; weakly informative priors are a
  more flexible default in modern workflows.
\item
  Jeffreys priors can behave badly in hierarchical models --- the
  multivariate version often produces improper posteriors when variance
  components are at the boundary. Reach for weakly informative
  half-Cauchy or half-Normal priors instead.
\item
  Always check that the resulting posterior is proper, especially when
  the prior itself is improper; a back-of-the-envelope check is to
  confirm that the posterior integrates to a finite quantity.
\item
  Reference priors (Bernardo) generalise Jeffreys' rule for ordered
  parameters of interest and nuisance components; they are the modern
  theoretical successor when an objective prior is genuinely required.
\item
  For pedagogical purposes, Jeffreys' prior is invaluable because it
  makes invariance under reparameterisation a hands-on exercise; for
  applied analyses, prefer weakly informative defaults unless invariance
  is a publication requirement.
\end{itemize}

\subsection{R Packages Used}\label{r-packages-used}

Base R suffices for closed-form Jeffreys updates in conjugate models;
\texttt{BayesianTools} and \texttt{mcmc} allow Jeffreys priors via
custom log-prior functions; \texttt{bayestestR::point\_estimate()} and
\texttt{HDInterval::hdi()} deliver tidy summaries on the resulting
posteriors.

\subsection{For Reviewers}\label{for-reviewers}

\textbf{What to look for in a paper using this method.}

\begin{itemize}
\tightlist
\item
  \textbf{Common misapplications.}
\item
  \textbf{Diagnostics that should be reported but often aren't.}
\item
  \textbf{Red flags in tables and figures.}
\item
  \textbf{What to verify.}
\item
  \textbf{What an adequate Methods paragraph must contain.}
\end{itemize}

\subsection{See also --- labs in this
chapter}\label{see-also-labs-in-this-chapter}

\begin{itemize}
\tightlist
\item
  \href{labs/lab-bayesian-thinking-control-vs-decision-error.Rmd}{Bayesian
  thinking; α control vs decision error}
\item
  \href{labs/lab-brms-stan-loo-hierarchical-models.Rmd}{brms / Stan,
  LOO, hierarchical models}
\end{itemize}

\end{document}
